% Source: https://artofproblemsolving.com/community/c3h1614410_scrabbler94_mock_mathcounts_sprint_17
% Stable identity: aops-topic-1614410-q01
Let the product be $P=n^3-1=(n-1)\left(n^2+n+1\right)$.

\textbf{Size.} The product of four different primes below $20$ is at least $2\cdot3\cdot5\cdot7=210$ and at most $11\cdot13\cdot17\cdot19=46189$. So $210\le n^3-1\le46189$, i.e. $6\le n\le35$. The product is squarefree, so both factors $n-1$ and $n^2+n+1$ are squarefree and have no common prime.

\textbf{Primes dividing $n^2+n+1$.} This number is odd. If a prime $p\ne3$ divides it, then $n^3\equiv1\pmod p$ but $n\not\equiv1$ (otherwise $p\mid3$), so $n$ has order $3$ modulo $p$ and $3\mid p-1$. Among the primes below $20$, only $3,7,13,19$ are $3$ or $\equiv1\pmod3$. So $n^2+n+1$ is a product of distinct primes from $\{3,7,13,19\}$, and for $6\le n\le35$ it lies between $43$ and $1261$.

The products of distinct primes from this set in that range are $57,91,133,247,273,399,741$. Solving $n^2+n+1=m$ for these, only $57$ ($n=7$), $91$ ($n=9$), $133$ ($n=11$) and $273$ ($n=16$) give integers; for instance $15^2+15+1=241<247<16^2+16+1$.

\textbf{Checking the candidates.}
\begin{itemize}
\item $n=7$: $6\cdot57=2\cdot3\cdot3\cdot19$, repeated $3$.
\item $n=9$: $8\cdot91=2^3\cdot7\cdot13$, not squarefree.
\item $n=11$: $10\cdot133=2\cdot5\cdot7\cdot19=1330=11^3-1$. This works.
\item $n=16$: $15\cdot273=3\cdot5\cdot3\cdot7\cdot13$, repeated $3$ (and five primes).
\end{itemize}
The primes are $2,5,7,19$, with sum $\boxed{33}$.
